% GENERATED FILE. Edit canonical lessons, then run npm run generate:books.
% canonical-source-hash: fnv1a64:8908f1ec9c14520f
% canonical-lessons: TA-C484-potuku, TA-C484-pen, TA-C484-valukkai, TA-C484-utir, TA-C484-narai, TA-R484-review-477-484

\chapter{Hair and Scalp}
\label{ch:ta-hair-and-scalp}
\hlchaptermodality{484}

\begin{chapteropening}
\textbf{By the end of this chapter:} \emph{I can talk about dandruff, lice, baldness and hair going grey.}

Complete the last lesson of chapter 484: I can talk about dandruff, lice, baldness and hair going grey.
\end{chapteropening}

\section[\texorpdfstring{\ta{பொடுகு} (poṭuku)}{பொடுகு (poṭuku)}]{\ta{பொடுகு} (poṭuku) — dandruff}

\label{lesson:TA-C484-potuku}



\begin{quote}
Before the new one: say the Tamil for to scratch an itch, then the Tamil for a mole.
\end{quote}

\subsection*{You'll want to know: \ta{பொடுகு}}
\textbf{\ta{பொடுகு}} — \emph{poṭuku} — \textquotedblleft{}dandruff\textquotedblright{}.

\begin{quote}
\ta{குளிர்} \ta{காலத்தில்} \ta{எனக்குப்} \ta{பொடுகு} \ta{அதிகமாகும்}.
\end{quote}

\begin{quote}
\emph{kuḷir kālattil eṉakkup poṭuku atikamākum.}
\end{quote}

\begin{quote}
In winter my dandruff gets worse.
\end{quote}

A mass noun: no plural.

\begin{cousinweb}
The origin of \textbf{\ta{பொடுகு}} (\emph{poṭuku}) is not settled; we could not find it in DEDR.
\end{cousinweb}

\begin{grammarlens}[title={What you've built}]
One more word for describing what you lived and felt.
\end{grammarlens}

\subsection*{Your turn}
\begin{itemize}
\raggedright
  \item \emph{Say it:} \emph{poṭuku}
  \item \emph{Say it:} \emph{poṭuku}, once more
  \item \emph{Say it:} \emph{kuḷir kālattil eṉakkup poṭuku atikamākum.} — In winter my dandruff gets worse.
  \item \emph{Recall:} say the Tamil for to scratch an itch, then the Tamil for a mole, then say \emph{poṭuku} again
\end{itemize}

\begin{tcolorbox}[breakable,colback=teal!4,colframe=teal!35!black,title={Before you move on}]
What does \ta{பொடுகு} mean? (\textquotedblleft{}Dandruff\textquotedblright{}.) Say it once more.
\end{tcolorbox}
\section[\texorpdfstring{\ta{பேன்} (pēṉ)}{பேன் (pēṉ)}]{\ta{பேன்} (pēṉ) — a louse}

\label{lesson:TA-C484-pen}



\begin{quote}
Before the new one: say the Tamil for a mole, then the Tamil for dandruff.
\end{quote}

\subsection*{You'll want to know: \ta{பேன்}}
\textbf{\ta{பேன்}} — \emph{pēṉ} — \textquotedblleft{}a louse\textquotedblright{}.

\begin{quote}
\ta{பள்ளியிலிருந்து} \ta{வந்த} \ta{குழந்தையின்} \ta{தலையில்} \ta{பேன்} \ta{இருந்தது}.
\end{quote}

\begin{quote}
\emph{paḷḷiyiliruntu vanta kuḻantaiyiṉ talaiyil pēṉ iruntatu.}
\end{quote}

\begin{quote}
The child back from school had lice in her hair.
\end{quote}

Plural: \ta{பேன்கள்}, \emph{pēṉkaḷ}.

Plural \ta{பேன்கள்}.

\begin{cousinweb}
\textbf{\ta{பேன்}} (\emph{pēṉ}) is native Dravidian (DEDR 4449); Malayalam, Kota and Kannada have the same word.
\end{cousinweb}

\begin{grammarlens}[title={What you've built}]
One more word for describing what you lived and felt.
\end{grammarlens}

\subsection*{Your turn}
\begin{itemize}
\raggedright
  \item \emph{Say it:} \emph{pēṉ}
  \item \emph{Say it:} \emph{pēṉ}, once more
  \item \emph{Say it:} \emph{paḷḷiyiliruntu vanta kuḻantaiyiṉ talaiyil pēṉ iruntatu.} — The child back from school had lice in her hair.
  \item \emph{Recall:} say the Tamil for a mole, then the Tamil for dandruff, then say \emph{pēṉ} again
\end{itemize}

\begin{tcolorbox}[breakable,colback=teal!4,colframe=teal!35!black,title={Before you move on}]
What does \ta{பேன்} mean? (\textquotedblleft{}A louse\textquotedblright{}.) Say it once more.
\end{tcolorbox}
\section[\texorpdfstring{\ta{வழுக்கை} (vaḻukkai)}{வழுக்கை (vaḻukkai)}]{\ta{வழுக்கை} (vaḻukkai) — baldness; a bald patch}

\label{lesson:TA-C484-valukkai}



\begin{quote}
Before the new one: say the Tamil for dandruff, then the Tamil for a louse.
\end{quote}

\subsection*{You'll want to know: \ta{வழுக்கை}}
\textbf{\ta{வழுக்கை}} — \emph{vaḻukkai} — \textquotedblleft{}baldness; a bald patch\textquotedblright{}.

\begin{quote}
\ta{நாற்பது} \ta{வயதில்} \ta{அப்பாவுக்கு} \ta{வழுக்கை} \ta{விழுந்தது}.
\end{quote}

\begin{quote}
\emph{nāṟpatu vayatil appāvukku vaḻukkai viḻuntatu.}
\end{quote}

\begin{quote}
Father went bald at forty.
\end{quote}

\ta{வழுக்கை} \ta{விழு}, to go bald.

\begin{cousinweb}
The origin of \textbf{\ta{வழுக்கை}} (\emph{vaḻukkai}) in this sense is not settled; DEDR lists the word only for the pulp of a tender coconut.
\end{cousinweb}

\begin{grammarlens}[title={What you've built}]
One more word for describing what you lived and felt.
\end{grammarlens}

\subsection*{Your turn}
\begin{itemize}
\raggedright
  \item \emph{Say it:} \emph{vaḻukkai}
  \item \emph{Say it:} \emph{vaḻukkai}, once more
  \item \emph{Say it:} \emph{nāṟpatu vayatil appāvukku vaḻukkai viḻuntatu.} — Father went bald at forty.
  \item \emph{Recall:} say the Tamil for dandruff, then the Tamil for a louse, then say \emph{vaḻukkai} again
\end{itemize}

\begin{tcolorbox}[breakable,colback=teal!4,colframe=teal!35!black,title={Before you move on}]
What does \ta{வழுக்கை} mean? (\textquotedblleft{}Baldness; a bald patch\textquotedblright{}.) Say it once more.
\end{tcolorbox}
\section[\texorpdfstring{\ta{உதிர்} (utir)}{உதிர் (utir)}]{\ta{உதிர்} (utir) — to fall out; to drop off}

\label{lesson:TA-C484-utir}



\begin{quote}
Before the new one: say the Tamil for a louse, then the Tamil for baldness.
\end{quote}

\subsection*{You'll want to know: \ta{உதிர்}}
\textbf{\ta{உதிர்}} — \emph{utir} — \textquotedblleft{}to fall out; to drop off\textquotedblright{}.

\begin{quote}
\ta{கவலையால்} \ta{அவருடைய} \ta{முடி} \ta{உதிர்ந்தது}.
\end{quote}

\begin{quote}
\emph{kavalaiyāl avaruṭaiya muṭi utirntatu.}
\end{quote}

\begin{quote}
His hair fell out from worry.
\end{quote}

Past (it …): \ta{உதிர்ந்தது}, \emph{utirntatu}.

Of hair, leaves and petals: the subject is the thing that falls, so the past is \ta{உதிர்ந்தது}.

\begin{cousinweb}
\textbf{\ta{உதிர்}} (\emph{utir}) is native Dravidian (DEDR 615), for leaves, fruit, hair and tears.
\end{cousinweb}

\begin{grammarlens}[title={What you've built}]
One more word for describing what you lived and felt.
\end{grammarlens}

\subsection*{Your turn}
\begin{itemize}
\raggedright
  \item \emph{Say it:} \emph{utir}
  \item \emph{Say it:} \emph{utir}, once more
  \item \emph{Say it:} \emph{kavalaiyāl avaruṭaiya muṭi utirntatu.} — His hair fell out from worry.
  \item \emph{Recall:} say the Tamil for a louse, then the Tamil for baldness, then say \emph{utir} again
\end{itemize}

\begin{tcolorbox}[breakable,colback=teal!4,colframe=teal!35!black,title={Before you move on}]
What does \ta{உதிர்} mean? (\textquotedblleft{}To fall out; to drop off\textquotedblright{}.) Say it once more.
\end{tcolorbox}
\section[\texorpdfstring{\ta{நரை} (narai)}{நரை (narai)}]{\ta{நரை} (narai) — to turn grey}

\label{lesson:TA-C484-narai}



\begin{quote}
Before the new one: say the Tamil for baldness, then the Tamil for to fall out.
\end{quote}

\subsection*{You'll want to know: \ta{நரை}}
\textbf{\ta{நரை}} — \emph{narai} — \textquotedblleft{}to turn grey\textquotedblright{}.

\begin{quote}
\ta{ஐம்பது} \ta{வயதில்} \ta{அம்மாவின்} \ta{முடி} \ta{நரைத்தது}.
\end{quote}

\begin{quote}
\emph{aimpatu vayatil ammāviṉ muṭi naraittatu.}
\end{quote}

\begin{quote}
At fifty mother’s hair turned grey.
\end{quote}

Past (it …): \ta{நரைத்தது}, \emph{naraittatu}.

As a noun, grey hair: \ta{நரை} \ta{முடி}.

\begin{cousinweb}
\textbf{\ta{நரை}} (\emph{narai}) is native Dravidian (DEDR 3609), both the grey hair and the greying.
\end{cousinweb}

\begin{grammarlens}[title={What you've built}]
One more word for describing what you lived and felt.
\end{grammarlens}

\subsection*{Your turn}
\begin{itemize}
\raggedright
  \item \emph{Say it:} \emph{narai}
  \item \emph{Say it:} \emph{narai}, once more
  \item \emph{Say it:} \emph{aimpatu vayatil ammāviṉ muṭi naraittatu.} — At fifty mother’s hair turned grey.
  \item \emph{Recall:} say the Tamil for baldness, then the Tamil for to fall out, then say \emph{narai} again
\end{itemize}

\begin{tcolorbox}[breakable,colback=teal!4,colframe=teal!35!black,title={Before you move on}]
What does \ta{நரை} mean? (\textquotedblleft{}To turn grey\textquotedblright{}.) Say it once more.
\end{tcolorbox}
\section[(dialogue)]{Looking back: from Serious Illness to Hair and Scalp}

\label{lesson:TA-R484-review-477-484}



\begin{quote}
Say the Tamil for a heart attack and cancer.
\end{quote}

\subsection*{Your turn}
Say these aloud:

\begin{itemize}
\raggedright
  \item a heart attack, cancer, a stroke, asthma, jaundice
  \item an organ, a lung, the liver, a kidney, the intestines
  \item a muscle, a nerve, the bosom, a rib, to throb
  \item the skull, the jaw, an eyebrow, an eyelid, to blink
  \item a thigh, the heel, to limp, to trip, to falter
  \item a sprain, to be dislocated, to get hurt, to snap, to ooze
  \item a sore, a blister, a pimple, to scratch an itch, a mole
  \item dandruff, a louse, baldness, to fall out, to turn grey
\end{itemize}

\begin{tcolorbox}[breakable,colback=teal!4,colframe=teal!35!black,title={Before you move on}]
A heart attack? (\textbf{\ta{மாரடைப்பு}}, \emph{māraṭaippu}.) Cancer? (\textbf{\ta{புற்றுநோய்}}, \emph{puṟṟunōy}.) A stroke? (\textbf{\ta{பக்கவாதம்}}, \emph{pakkavātam}.) Asthma? (\textbf{\ta{ஆஸ்துமா}}, \emph{āstumā}.) Jaundice? (\textbf{\ta{காமாலை}}, \emph{kāmālai}.) An organ? (\textbf{\ta{உறுப்பு}}, \emph{uṟuppu}.) A lung? (\textbf{\ta{நுரையீரல்}}, \emph{nuraiyīral}.) The liver? (\textbf{\ta{கல்லீரல்}}, \emph{kallīral}.) A kidney? (\textbf{\ta{சிறுநீரகம்}}, \emph{ciṟunīrakam}.) The intestines? (\textbf{\ta{குடல்}}, \emph{kuṭal}.) A muscle? (\textbf{\ta{தசை}}, \emph{tacai}.) A nerve? (\textbf{\ta{நரம்பு}}, \emph{narampu}.) The bosom? (\textbf{\ta{நெஞ்சு}}, \emph{neñcu}.) A rib? (\textbf{\ta{விலா}}, \emph{vilā}.) To throb? (\textbf{\ta{துடி}}, \emph{tuṭi}.) The skull? (\textbf{\ta{மண்டை}}, \emph{maṇṭai}.) The jaw? (\textbf{\ta{தாடை}}, \emph{tāṭai}.) An eyebrow? (\textbf{\ta{புருவம்}}, \emph{puruvam}.) An eyelid? (\textbf{\ta{இமை}}, \emph{imai}.) To blink? (\textbf{\ta{சிமிட்டு}}, \emph{cimiṭṭu}.) A thigh? (\textbf{\ta{தொடை}}, \emph{toṭai}.) The heel? (\textbf{\ta{குதிகால்}}, \emph{kutikāl}.) To limp? (\textbf{\ta{நொண்டு}}, \emph{noṇṭu}.) To trip? (\textbf{\ta{இடறு}}, \emph{iṭaṟu}.) To falter? (\textbf{\ta{தடுமாறு}}, \emph{taṭumāṟu}.) A sprain? (\textbf{\ta{சுளுக்கு}}, \emph{cuḷukku}.) To be dislocated? (\textbf{\ta{பிசகு}}, \emph{picaku}.) To get hurt? (\textbf{\ta{அடிபடு}}, \emph{aṭipaṭu}.) To snap? (\textbf{\ta{முறி}}, \emph{muṟi}.) To ooze? (\textbf{\ta{கசி}}, \emph{kaci}.) A sore? (\textbf{\ta{புண்}}, \emph{puṇ}.) A blister? (\textbf{\ta{கொப்புளம்}}, \emph{koppuḷam}.) A pimple? (\textbf{\ta{பரு}}, \emph{paru}.) To scratch an itch? (\textbf{\ta{சொறி}}, \emph{coṟi}.) A mole? (\textbf{\ta{மச்சம்}}, \emph{maccam}.) Dandruff? (\textbf{\ta{பொடுகு}}, \emph{poṭuku}.) A louse? (\textbf{\ta{பேன்}}, \emph{pēṉ}.) Baldness? (\textbf{\ta{வழுக்கை}}, \emph{vaḻukkai}.) To fall out? (\textbf{\ta{உதிர்}}, \emph{utir}.) To turn grey? (\textbf{\ta{நரை}}, \emph{narai}.)
\end{tcolorbox}
